\section{Modules}

\begin{exercise}\label{ex2.1}
	B\'{e}zout's identity.
\end{exercise}

\begin{exercise}\label{ex2.2}
	Tensor the exact sequence $0\rightarrow\mathfrak{a}\rightarrow A\rightarrow A/\mathfrak{a}\rightarrow0$. Calculate the image of  $\mathfrak{a}\otimes M\rightarrow A\otimes M$.
\end{exercise}

\begin{exercise}\label{ex2.3}
	Consider the residue field $k$ of $A$ and pass $M\otimes_AN$ to vector space $M_k\otimes_kN_k$. Use \Cref{ex2.2} and \hyperref[2.6]{Nakayama's lemma}.
\end{exercise}

\begin{exercise}\label{ex2.4}
	There exist canonical projections.
\end{exercise}

\begin{exercise}\label{ex2.5}
	The definition of flatness only considers the module structure. Use \Cref{ex2.4}.
\end{exercise}

\begin{exercise}\label{ex2.6}
	Construct two homomorphisms.
\end{exercise}

\begin{exercise}\label{ex2.7}
	Proof that $(A/\mathfrak{p})[x]$ and $A[x]/\mathfrak{p}[x]$ are isomorphic.
\end{exercise}

\begin{exercise}\label{ex2.8}
	\begin{enumerate}
		\item[(ii)] Use \Cref{2.15}.
	\end{enumerate}
\end{exercise}

\begin{exercise}\label{ex2.9}
	Pick $y_1,\cdots,y_m\in M$ whose images are the generators of $M^{\prime\prime}$. Use diagram chasing for an element of $M$. 
\end{exercise}

\begin{exercise}\label{ex2.10}
	$u(m)+\mathfrak{a}N\in N$, then use surjectiveness. Apply \hyperref[2.6]{Nakayama's lemma} on $(u(M)+\mathfrak{a}N)/N$.
\end{exercise}

\begin{exercise}\label{ex2.11}
	\begin{enumerate}
		\item[(i),(ii)] Consider a residue field and pass to vector spaces. Use the right exactness of tensor product.
		\item[(iii)] If $m>n$, let $\psi:A^m\xrightarrow{\phi}A^n\hookrightarrow A^m$. By \Cref{2.4}, $\psi$ satisfies $\psi^n+a_{n-1}\psi^{n-1}+\cdots+a_0=0$. Pick such equation to have the minimum degree, then $\psi$ injective implies $a_0\neq0$. Apply the equation to element $(0,\cdots,0,1)$ to obtain contradiction.
	\end{enumerate}
\end{exercise}

\begin{exercise}\label{ex2.12}
	Choose $u_i\in M$ such that $\phi\left(u_i\right)$ are a basis of $A^n$, then proof that $0\rightarrow\ker(\phi)\rightarrow M\xrightarrow{\phi}A^n\rightarrow0$ splits.
\end{exercise}

\begin{exercise}\label{ex2.13}
	$p\circ g=\operatorname{id}_N$ injective $\Rightarrow$ $g$ injective, so $0\rightarrow N\xrightarrow{g}N_B\rightarrow N_B/g(N)\rightarrow0$ splits.
\end{exercise}

\begin{exercise}\label{ex2.14}
	This question asks to proof $\mu_i=\mu_j\circ\mu_{ij}$.
\end{exercise}

\begin{exercise}\label{ex2.15}
	\begin{enumerate}[label=(\roman*),ref=\theexercise(\roman*)]\crefalias{enumi}{exercise}
		\item\label{ex2.15.1} Use \Cref{ex2.14} to increase indices.
		\item\label{ex2.15.2} \textcolor{magenta}{$x_i=\sum_k\left(y_k-\mu_{kl}\left(y_k\right)\right)$ is an equation in $\bigoplus M_i$, so $\mu_{ij}$ cannot be applied directly.} Use canonical projection $\pi_i:\bigoplus M_i\rightarrow M_i$ on the equation first.
	\end{enumerate}
\end{exercise}

\begin{exercise}[Universal property of direct limit]\label{ex2.16}
	Use \Cref{ex2.15.1} for both existence and uniqueness.
\end{exercise}

\begin{exercise}\label{ex2.17}
	Similar to \Cref{ex2.15.1}, move components into a same index.
\end{exercise}

\begin{exercise}\label{ex2.18}
	Use \Cref{ex2.16}.
\end{exercise}

\begin{exercise}[$\varinjlim$ is an exact functor]\label{ex2.19}
	Diagram chasing. Use \Cref{ex2.15}.
\end{exercise}

\begin{exercise}\label{ex2.20}
	Draw the diagram. Define two homomorphisms by the universal properties of tensor product and direct limit. Verify their compositions are identity maps. 
\end{exercise}

\begin{exercise}\label{ex2.21}
	Use \Cref{ex2.15.2} on $\alpha_i(1)$.
\end{exercise}

\begin{exercise}\label{ex2.22}
	Use \Cref{ex2.15}. Pick a nilpotent or assume $xy=0$.
\end{exercise}

\begin{exercise}\label{ex2.23}
	If $J\subseteq J^\prime$, $B_J\cong B_J\otimes\left(A\right)_{i\in J^\prime\backslash J}$. Let $\Gamma$ be the set of all finite subsets of $\Lambda$, then $\left(B_J\right)_{J\in\Gamma}$ is the case of \Cref{ex2.21}. We can also use \Cref{ex2.17} to see its form.
\end{exercise}

\begin{exercise}\label{ex2.24}
	Directly from the definition of $\operatorname{Tor}$ functor.
\end{exercise}

\begin{exercise}\label{ex2.25}
	Use \hyperref[2.10]{Snake lemma}.
\end{exercise}

\begin{exercise}\label{ex2.26}
	Let $M_n$ be generated by $\left\{x_1,\cdots,x_n\right\}$. We have an exact sequence
	\begin{equation*}
		\operatorname{Tor}_1\left(M_{n-1},N\right)\rightarrow\operatorname{Tor}_1\left(M_n,N\right)\rightarrow\operatorname{Tor}_1\left(M_n/M_{n-1},N\right).
	\end{equation*}
	By \Cref{2.3}, $M_n/M_{n-1}\cong A/\mathfrak{a}$.
	
	Let $f:\mathfrak{a}\otimes N\rightarrow A\otimes N$, $f\left(\sum a_i\otimes y_i\right)=0$. Let $\mathfrak{a}^\prime$ be generated by $a_i$. By assumption, $f|_{\mathfrak{a}^\prime\otimes N}$ is injective, so $\sum a_i\otimes y_i=0$ and thus $\operatorname{Tor}_1(A/\mathfrak{a},N)=0$. Since $M_0=A$ and $\operatorname{Tor}_1\left(A,N\right)=0$, the exact sequence implies $\operatorname{Tor}_1\left(M_n,N\right)=0$ for all $n$. Then use \Cref{2.19.4}.
	
	\textcolor{green!70!black}{Another approach is to use the fact that any module is the direct limit of some finitely generated modules, and direct limit is an exact functor.}
\end{exercise}

\begin{exercise}\label{ex2.27}
		\begin{enumerate}
		\item[(i) $\Rightarrow$ (ii)] $A/(x)$ is a flat $A$-module, then $\phi:A/(x)\otimes(x)\rightarrow A/(x)$ injective. $\phi(1\otimes x)=x=0$, so $A/(x)\otimes(x)=0$. Then use \Cref{ex2.2}.
		\item[(ii) $\Rightarrow$ (iii)] Let $x=ax^2$, then $ax$ is an idempotent. $(x)\ni bx=bax^2=(bx)(ax)\in(ax)$, so $(ax)=(x)$. Similar to \Cref{ex1.11.3}, every finitely generated ideal is generated by an idempotent $e$, so $A=(e)\oplus(1-e)$.
		\item[(iii) $\Rightarrow$ (i)] $N\otimes_AA=\left(N\otimes_A\mathfrak{a}\right)\oplus\left(N\otimes_A\mathfrak{a}^\prime\right)$, so $\left(N\otimes_A\mathfrak{a}\right)\rightarrow N\otimes_AA$ injective. Use \Cref{ex2.26}.
	\end{enumerate}
\end{exercise}

\begin{exercise}\label{ex2.28}
	Use \hyperref[ex2.27]{\Cref*{ex2.27}(ii)}.
\end{exercise}